% ═══════════════════════════════════════════════════════════════════
%  Section 7 — Discussion
% ═══════════════════════════════════════════════════════════════════

\label{sec:discussion}

This section interprets the experimental results, examines whether LLM agents produce behaviourally recognisable negotiation patterns, identifies the principal methodological tensions, situates the findings relative to prior work, and identifies the limitations of the current study.

% ───────────────────────────────────────────────────────────────────
\subsection{Behavioural Findings: What the Experiments Reveal About LLM Negotiation}
\label{sec:disc-positive}

The experimental programme reveals a consistent set of behavioural patterns that emerge from LLM-driven negotiation without gate-driven enforcement (all acceptances across the programme are LLM-driven; prompt steering remains active as an uncontrolled variable).

\paragraph{Concession and bargaining dynamics.}
LLM agents produce concave (Boulware-style) concession trajectories (Finding~1): the buyer makes a large opening concession of \$13.31 between rounds~0 and~2, then hardens rapidly.  The seller exhibits strong anchor rigidity, barely moving from its initial position (Finding~3).  The resulting 5.65:1 buyer-to-seller concession ratio indicates an asymmetric bargaining dynamic in which the buyer bears the concession burden.  Deadline effects are clearly visible: under the 6-round horizon, 82.4\% of all deals close in the final two rounds, with a 25\% higher concession rate compared to the 20-round horizon (Finding~4).

\paragraph{Deal success and market dynamics.}
Deal success rates vary systematically from 97.2\% in fixed-scenario experiments to 66.7--72.4\% in distribution-mode market experiments, primarily reflecting ZOPA availability.  Market liquidity of 81\% (Experiment~3) indicates that the Qwen2.5-14B model is effective at reaching deals within 10 rounds without gate assistance.  Prices exhibit persistent within-tick dispersion without convergence (Experiment~3), and directionally correct but substantially attenuated responses to supply and demand shocks (Experiments~4 and~6).

\paragraph{Informational anchor sensitivity.}
the Experiment~4 extension demonstrates that the environmental reference price in the agent prompt is a primary determinant of price levels: mean transaction prices span approximately \$30 across anchor modes under identical fundamentals.  This sensitivity extends to shock responsiveness, with the shock effect sign reversing between fixed-anchor and updated-anchor modes.  This finding establishes that LLM agents do not independently discover prices from private constraints alone; the informational environment materially shapes bargaining outcomes.

\paragraph{Micro--macro capability gap.}
A central insight running across the experimental programme is that LLM agents demonstrate separable capabilities in deal-closing and price-discovery.  Deal-closing performance is high: market liquidity reaches 81\% without gate assistance, and individual sessions settle reliably within the round limit.  Price-discovery performance is low: price elasticities to fundamental shocks are 0.035--0.275, and prices track the prompt-level reference signal rather than the agents' private constraints.  This dissociation---strong convergence to agreement, weak convergence to informationally correct prices---is the primary structural finding of the project and organises the discussion in Sections~\ref{sec:disc-failure} and~\ref{sec:disc-micro-macro}.

\paragraph{Institutional mechanism effects.}
Surplus-maximising matching improves total welfare by 7.2\% relative to random matching (Experiment~5), with gains shared across both buyers and sellers.  This demonstrates that mechanism design can improve market outcomes independently of agent sophistication.

% ───────────────────────────────────────────────────────────────────
\subsection{LLM Agents and Human Negotiation Behaviour}
\label{sec:disc-human}

A central question motivating this project is whether LLM agents produce negotiation behaviour that is recognisably similar to patterns documented in human negotiation studies.  The experimental programme enables directional comparison against three well-established human bargaining patterns, though all comparisons are indirect: no human participants were tested under the same protocol.

\paragraph{Where patterns align.}
Deadline-driven concession acceleration (Experiment~2) is directionally consistent with Roth et al.'s~\cite{roth1988deadline} documentation of deadline effects in human bargaining: agents make disproportionately large concessions near the deadline under time pressure.  The concave concession trajectory (Experiment~1) is consistent with the Boulware bargaining pattern~\cite{raiffa1982} documented in the negotiation literature, in which a negotiator makes early accommodations but becomes increasingly resistant.  Reference-price sensitivity (the Experiment~4 extension) parallels the first-offer anchoring effect~\cite{galinsky2001}, though through a different mechanism: environmental prompt-level anchoring rather than opponent-generated first offers.

\paragraph{Where patterns diverge.}
Prices do not shift meaningfully with deadlines (Experiment~2, Finding~5): mean deal prices differ by less than \$2 across horizon conditions, whereas human deadline studies typically document both timing and price effects.  This dissociation between timing sensitivity and price sensitivity under deadline conditions is a specific, testable behavioural difference that distinguishes LLM negotiation from human negotiation.  In human studies, deadline pressure typically induces larger concessions---affecting both when and at what price deals close; LLM agents show only the former effect, suggesting they respond to deadline salience as a procedural cue to close the deal rather than as strategic pressure to revise their price target.  Under the finite-horizon counteroffer protocol, agents do not voluntarily walk away from profitable sessions; the round limit serves as the outside option rather than an active rejection decision.  This protocol eliminates the threat-of-rejection dynamic that is central to strategic human negotiation behaviour.  Price stickiness (Experiment~6) indicates that agents do not efficiently update prices in response to changing fundamentals, a form of departure from the adaptive price-discovery behaviour expected of experienced human negotiators.

\paragraph{Conservative interpretation.}
The observed alignments are consistent with LLM agents having absorbed negotiation-relevant behavioural patterns from training data, but this cannot be confirmed from the current experiments.  The appropriate framing is that LLM agents exhibit qualitative behavioural patterns that are interpretable through the same theoretical lens as human negotiation behaviour, without claiming that the underlying cognitive or strategic mechanisms are similar.  Direct human comparison under identical protocol would be required to substantiate stronger claims.

% ───────────────────────────────────────────────────────────────────
\subsection{Price Stickiness and Weak Price Discovery}
\label{sec:disc-failure}

A substantively important finding is the systematic failure of LLM agents to pass through fundamental changes into prices.  In Experiment~6, a \$20 shift in buyer value produces only a \$4.63 price increase (elasticity 0.275), and a \$20 shift in seller cost produces only a \$1.17 price increase (elasticity 0.035).  Agents adjust prices only weakly in response to changing fundamentals, even when those fundamentals are encoded in their private constraints.  the Experiment~4 extension identifies the prompt-level reference price as a substantial contributor, though the full mechanism is not isolated.  This price stickiness is observed in sessions where all settlements are LLM-driven and the gate is absent.  The parallel finding in Experiment~3---persistent within-tick price dispersion without convergence over 10 ticks---reinforces the same conclusion: under the tested model--prompt combination, agents function as effective deal-closing mechanisms but poor price-discovery mechanisms.

the Experiment~4 extension (Section~\ref{sec:results-anchor}) provides supporting evidence that the prompt-level reference price is one contributor to price stickiness.  Across the full range of anchor modes, mean transaction prices vary by approximately \$30 (from \$79 under \texttt{updated\_anchor} to \$109 under \texttt{fixed\_anchor}), indicating that the environmental reference signal influences price levels substantially.  A post-hoc pooled comparison across 1{,}388 deals (Section~\ref{sec:results-anchor}) confirms that all three pairwise price differences are statistically distinguishable after Bonferroni correction, with effect sizes ranging from $d = 0.48$ (fixed vs.\ no-anchor) to $d = 1.32$ (fixed vs.\ updated), indicating that anchor-mode effects are practically as well as statistically meaningful under this protocol.  The \texttt{no\_anchor} condition also eliminates shock sensitivity almost entirely (shock price effect of $-$\$0.20), indicating that without an external reference signal, agents default to an internal prior and do not adjust prices in response to fundamental changes.  The \texttt{updated\_anchor} condition, which feeds prior deal prices back into the prompt, produces a self-reinforcing downward drift and reverses the sign of the shock effect.  The pooled shock-sensitivity analysis additionally shows that the updated-anchor condition is the only one that produces a statistically significant mean price difference between shock and no-shock sessions ($p = 0.032$, $d = 0.21$); fixed-anchor and no-anchor conditions show no significant pooled shock effect ($p = 0.77$ and $p = 0.92$ respectively).  This result is consistent with Findings~13 and~14 but should be interpreted cautiously: the updated-anchor shock effect is marginal ($d = 0.21$, a small effect by conventional benchmarks) and the analysis is post-hoc with sessions clustered within seeds.  These findings suggest that price stickiness in this simulator is not solely a property of the LLM's reasoning limitations but is partly a consequence of the environmental information structure: agents are sensitive to the reference the prompt provides.  The full mechanism underlying price stickiness---whether it arises primarily from training-data priors, prompt-level anchoring, or the conservative opening strategy observed in Experiment~1---is not resolved by the current experiments.  This limitation also has a practical dimension: deployed LLM negotiation systems using similar prompt architectures may exhibit systematic mis-pricing that does not self-correct as market conditions change, because agents anchor to the reference price embedded in the prompt rather than adapting to the informational content of offers received.

Price stickiness also has an important implication for interpreting welfare results.  As noted in Section~\ref{sec:results-supply-demand}, total welfare in this simulator is defined as $(v_b - c_s)$, which is independent of the deal price $p$.  Welfare comparisons across supply-demand conditions therefore reflect shifts in the composition of trading pairs, not agent-level price adjustment.  Welfare results look strong because they are partly mechanical; price results look weak because they require genuine LLM behavioural adaptation.  This distinction is the most important interpretive caution for Experiments~6 and~5.

\paragraph{Protocol design choices.}
The bargaining protocol implements a finite-horizon counteroffer format: agents counter-offer or accept, and episodes that reach the round limit end as no-deal outcomes.  The prompt additionally instructs agents to concede monotonically in the correct direction (buyers upward, sellers downward).  Both choices were introduced for practical reasons: without structured deal finalization, small language models frequently terminate profitable negotiations prematurely, producing near-zero deal rates; without directional concession guidance, agents sometimes concede in the wrong direction.  The concession \emph{magnitude} and \emph{timing} remain interpretable (the instruction constrains direction, not step size), but the monotonic concession pattern reported in Experiment~1 cannot be fully attributed to emergent LLM strategy.  This illustrates a recurring tension between engineering reliability and experimental cleanliness.

\paragraph{Demand-side and supply-side asymmetry.}
Experiment~6 identifies an asymmetry in the direction opposite to what classical theory predicts: the demand shock raises prices by \$4.63 (+4.2\%) while the supply shock raises them by only \$1.17 (+1.1\%), despite identical \$20 input magnitudes.  Standard supply-demand intuition suggests that a hard seller cost floor should produce stronger pass-through than a buyer value increase, but the observed result reverses this expectation.  A plausible explanation is that buyers, as first movers, anchor prices more strongly than sellers: LLM buyers open higher under a demand shock because their value ceiling rises, but LLM sellers do not commensurately raise their counter-offers under a supply shock.  Alternative explanations include the conservative buyer opening strategy documented in Experiment~1 (where free-language buyers open at \$71.88 against a \$120 value) and model-specific anchoring to a fixed reference range.  None of these explanations can be separated without targeted ablation.

The supply shock also reduces market liquidity from 0.724 to 0.667, consistent with a tighter ZOPA degrading LLM convergence even when agreement is structurally possible.

\paragraph{Prompt architecture as a methodological consideration.}
In human negotiation research, the cognitive architecture of the bargainer is relatively fixed; experimental variation comes from the task, the opponent, and the information structure.  In LLM negotiation, the prompt shapes the agent's decision process, and changes to prompt design---adding a reasoning scaffold, including or excluding surplus information, steering toward or away from rejection---can alter outcomes.  LLM negotiation results should therefore be understood as characterising a specific model--prompt--protocol combination rather than LLM capability in general.  The Experiment~1 comparison between rule-based and free-language agents shows that the decision architecture shapes offer behaviour: the free-language agent produces concave concession trajectories while the rule-based agent follows a fixed linear schedule.  Whether this reflects emergent strategic reasoning or prompt-induced output distribution cannot be determined from the current experiments.

% ───────────────────────────────────────────────────────────────────
\subsection{Micro--Macro Tension}
\label{sec:disc-micro-macro}

A structural tension runs through the experimental programme: designs that preserve raw agent behaviour may weaken aggregate market regularity, while stronger enforcement can stabilise outcomes but partially replaces agent decision-making.

As described in Section~\ref{sec:gate}, the simulator includes a feasibility gate that, when enabled, overrides agent decisions by automatically accepting feasible offers.  Preliminary runs with the gate enabled showed substantial gate involvement, producing near-100\% deal rates in wide-ZOPA conditions but simultaneously suppressing the very dynamics the experiments were designed to study.  With the gate disabled, two mediation layers remain active between the model's generation and the enacted action: the rethink loop corrects format and self-constraint violations inside the agent, and the ActionJudge enforces protocol and global-bound constraints at the session level.

With the gate off, agents produce concave concession trajectories (Finding~1) and deadline-sensitive timing (Finding~4)---patterns that preliminary gate-on runs showed were suppressed by enforcement.  Market liquidity is 81\% with zero gate involvement (Finding~7), indicating that the Qwen2.5-14B model can sustain a functioning market without enforcement assistance.  The micro--macro tension therefore manifests differently: enforcement is not needed to make the market function, but the absence of enforcement exposes a different limitation---price stickiness.

The price stickiness finding (elasticities 0.03--0.28, Finding~12) is the clearest example of this tension.  Individual agents close deals effectively (micro-level function), but they do not update their price targets efficiently in response to changing fundamentals.  The result at the market level is persistent price inertia: the market clears at prices that reflect the agents' anchored reference range rather than the shifted equilibrium implied by the new parameter distributions.  Effective deal closure (micro) and effective price discovery (macro) appear to be separable capabilities, with LLM agents demonstrating the former more reliably than the latter.

% ───────────────────────────────────────────────────────────────────
\subsection{Comparison to Prior Work}
\label{sec:disc-prior}

We compare our findings to the prior work reviewed in Section~\ref{sec:related}.

\paragraph{Generative agent simulations.}
Park et al.~\cite{park2023generative} demonstrated that LLM agents exhibit emergent social behaviours in open-ended environments.  Our market experiments (3 and~4) apply a related approach to an economic setting with hard structural constraints: LLM agents produce aggregate patterns---persistent price dispersion, varying liquidity, and surplus asymmetry---without being explicitly programmed to produce these outcomes.  The key difference is that our environment is constrained (hard budget and cost limits, structured protocol) rather than open-ended, enabling direct comparison against theoretical predictions.

\paragraph{LLM negotiation.}
Abdelnabi et al.~\cite{abdelnabi2024llmdeliberation} find that deliberation prompting improves agreement quality in multi-party, multi-issue negotiation games.  Our Experiment~1 provides a complementary finding in a simpler setting: the free-language agent produces concave concession trajectories with a 5.65:1 buyer-to-seller concession ratio, illustrating how decision architecture shapes bargaining dynamics even in a single-issue bilateral protocol.  Our setting is single-issue bilateral bargaining, which isolates price dynamics from the issue-trading effects present in their multi-issue design.

Mao et al.~\cite{mao2024personalities} show that personality-assigned prompts shift negotiation outcomes.  Our prompt architecture (acceptance hints, concession-direction instructions) functions similarly as a prompt-level modifier; the methodological difference is that we hold prompt architecture constant as a fixed protocol parameter rather than studying it as an independent variable.  All reported experiments use the default neutral communication strategy.

Lewis et al.~\cite{lewis2017dealornodeal} found that end-to-end trained negotiation agents develop non-human-like strategies, including degenerate repetitive behaviour.  Our LLM agents show a different form of departure from human norms: rather than degenerating, they fail to efficiently update prices in response to changing fundamentals while maintaining otherwise reasonable bargaining behaviour.  Both findings suggest that artificial agents do not spontaneously replicate the full spectrum of human negotiation behaviour, though the underlying mechanisms differ---reinforcement learning via self-play versus prompt-conditioned inference.

\paragraph{Classical theory.}
Several findings are directionally consistent with qualitative predictions from bilateral bargaining theory.  Market prices exhibit persistent within-tick dispersion consistent with independently negotiated bilateral pairs~\cite{rubinstein1985}, and the buyer-first protocol produces buyer-favourable surplus splits consistent with first-mover advantage~\cite{rubinstein1982}.  Experiment~6 shows that a +\$20 demand shock raises prices by \$4.63 (+4.2\%) and a +\$20 supply shock raises them by \$1.17 (+1.1\%), both directionally consistent with competitive market predictions but substantially attenuated.

The deadline result is directionally consistent with Roth et al.~\cite{roth1988deadline}: concession rates are 25\% higher under the 6-round horizon and 82.4\% of deals close in the final two rounds.  However, one finding runs counter to classical supply-demand theory: the demand shock produces a \$4.63 price increase (elasticity 0.275) while the symmetric supply shock produces only a \$1.17 increase (elasticity 0.035).  Standard theory predicts that a binding seller cost floor should produce at least as strong a price pass-through as a buyer value increase of the same magnitude; the observed result is the reverse.  This demand-supply price asymmetry is therefore a substantive negative finding, not merely an attenuated version of the classical prediction.

\paragraph{AgenticPay and open-weight model limitations.}
Liu et al.~\cite{liu2026agenticpay} introduce AgenticPay, a benchmark for multi-agent buyer--seller negotiation driven by natural language, spanning bilateral to many-to-many settings.  Their evaluation of open-weight models on 111 negotiation tasks provides an informative external reference for interpreting our results.  On the AgenticPay benchmark, Qwen3-14B---a model of the same nominal parameter scale (14B) as our \texttt{Qwen/Qwen2.5-14B-Instruct}, though from a different model generation---achieves a GlobalScore of 63.9 with a 20.7\% timeout rate, compared to 100\% deal rates and GlobalScores above 81 for frontier proprietary models.  Our own results show a 9.7\% timeout rate in the market dynamics experiment and price elasticities of 0.03--0.28, consistent with the AgenticPay finding that open-weight models engage in price discovery but struggle with long-horizon strategic convergence.  The AgenticPay benchmark provides an illustrative external reference, though the experimental tasks, protocols, and model generations differ substantially; no direct performance comparison between Qwen3-14B and Qwen2.5-14B-Instruct is intended.  Both studies observe that models of comparable but not identical specification handle the mechanics of bilateral bargaining while showing systematic limitations in price convergence, but the underlying mechanisms are not established by either study.

% ───────────────────────────────────────────────────────────────────
\subsection{Limitations}
\label{sec:disc-limitations}

We identify nine limitations of the current study, ordered by their impact on the interpretation of results.

\begin{enumerate}
    \item \textbf{Single model family; no cross-model comparison.}  All experiments use \texttt{Qwen/Qwen2.5-14B-Instruct} at temperature~0.2.  Results may not generalise to other model families, scales, or temperatures.  In particular, the emergence of concave concession curves and deadline effects may reflect model-specific properties rather than general properties of LLM bargaining.  No cross-model comparison is provided.

    \item \textbf{Price stickiness: finding and limitation simultaneously.}  The low price elasticities in Experiment~6 (0.035--0.275) are a substantively important finding, but they are also a limitation: agents do not discover prices efficiently, meaning the simulator cannot be used to test price discovery mechanisms that require agents to update prices in response to changing fundamentals.  The mechanism underlying stickiness---whether it arises from training-data anchoring, prompt-level reference framing, or the conservative opening strategy observed in Experiment~1---is not resolved by the current experiments.  At the same time, the finding itself is informative: it reveals that LLM agents are more responsive to environmental reference information embedded in the prompt than to private constraint information encoded in their own parameters, a distinction with implications for how such agents process strategic information.

    \item \textbf{No prompt ablation.}  The accept-condition block, monotonic concession instructions, the counteroffer-only action space, and deadline-salience language are all active across experiments but none are individually ablated.  Their separate effects on deal rates, concession shape, price stickiness, and deadline sensitivity are unknown.  The monotonic concession pattern (Finding~1) is partly attributable to the directional concession instruction; the seller's observed resistance to concession (Finding~3) may partly reflect prompt asymmetry rather than intrinsic role-based behaviour.

    \item \textbf{Three seeds is limited replication.}  All experiments use three seeds (42, 123, 456).  This provides basic cross-seed consistency checks and directional replication, but it is insufficient for hypothesis testing or effect-size estimation.  Seed-level variance is visible in the shock recovery experiment (Experiment~4), where individual seeds show divergent trajectories.  Most analyses are descriptive (means, standard deviations).  The post-hoc pooled anchor-mode comparison reported in Section~\ref{sec:results-anchor} constitutes a partial exception: it applies Welch's $t$-tests with Bonferroni correction to 1{,}388 pooled deal observations.  However, this analysis increases the number of observations rather than the number of independent seeds (still~3), and sessions within a seed share market context, meaning standard errors from session-level $t$-tests may be underestimated; the analysis should be treated as indicative rather than as a formally powered test.

    \item \textbf{No human baseline.}  All experiments compare LLM agents against rule-based or other LLM agents.  No human negotiation data is collected under the same protocol.  Claims about departure from or alignment with human negotiation behaviour therefore rely on the external literature rather than direct comparison.

    \item \textbf{Single-price negotiation.}  All sessions negotiate a single price.  Real markets involve quality, delivery terms, and bundling.  Multi-issue negotiation may produce qualitatively different dynamics, including the logrolling and issue-trading effects studied by Lewis et al.~\cite{lewis2017dealornodeal}.

    \item \textbf{No cross-run learning.}  Agents do not accumulate experience across sessions within the completed experiments.  The memory and reputation agents are implemented but untested.  Whether LLM agents adapt strategies over repeated interactions remains an open question.

    \item \textbf{LLM non-determinism.}  Even at temperature~0.2, model outputs are not fully deterministic.  The seeded RNG controls market-level randomness (population generation, matching, shocks) but not LLM inference randomness.  This introduces uncontrolled variance into all LLM-agent conditions.

    \item \textbf{Demand-supply price asymmetry mechanism unresolved.}  The demand shock produces a \$4.63 price increase (elasticity 0.275) while the supply shock produces only \$1.17 (elasticity 0.035) despite identical \$20 input magnitudes.  Three candidate explanations---buyer first-mover dominance, cost-floor anchoring by sellers, and nonlinear ZOPA-width sensitivity---are consistent with the data but have not been isolated through targeted ablation.
\end{enumerate}
